\section{Objectifs et aire d'étude}

\subsection{Objectifs}

Un avenant du projet R875.1 a été signé avec le Ministère en mai 2025. Cet avenant
révise le devis initial du projet de recherche R875.1 (V.2022-11-04) en raison du
report des travaux d'aménagement des passages fauniques du Ministère à une date
indéterminée. L'objectif principal du projet n'a pas changé, et consiste à produire
des connaissances sur les critères favorisant la connectivité écologique des
amphibiens, des reptiles, ainsi que des petits mammifères dans la région
métropolitaine de Montréal. Toutefois, les objectifs spécifiques et les hypothèses
ont été modifiés comme suit :

\begin{enumerate}
\item Développer des connaissances dans le domaine de l'écologie routière, en
        mettant l'accent sur les routes du Ministère.

\begin{description}
\item[Hypothèse 1.1 :] Les taux de mortalité des amphibiens diffèrent
                entre les espèces, mais coïncident avec les périodes de migration
                des adultes et de dispersion des juvéniles, et avec les périodes de
                précipitations.
\item[Hypothèse 1.2 :] Les taux de mortalité des tortues sont plus élevés
                aux tronçons de route adjacents à une combinaison de milieux humides
                et de milieux ouverts.
\item[Hypothèse 1.3 :] Les taux de mortalité des tortues et des couleuvres
                sont associés aux périodes de mouvement (ponte, dispersion de
                juvéniles).
\item[Hypothèse 1.4 :] Les taux de mortalité des petits mammifères sont
                plus élevés dans les tronçons adjacents à des milieux forestiers ou
                à des milieux humides.
                
\end{description}

\item Identifier les critères d'emplacements optimaux des passages fauniques par
        les rainettes faux-grillons de l'Ouest (RFGO) et les tortues (sélection des
        emplacements des structures, mortalité, connectivité) en fonction des types
        de routes et des saisons.

\begin{description}

\item[Hypothèse 2.1 :] L'emplacement optimal de futurs passages fauniques
                dépend de la qualité du milieu de part et d'autre de la route
                (milieu humide, milieu ouvert, milieu forestier), de l'intensité de
                la circulation sur la route et du taxon visé. Ces conditions
                rassemblées favorisent les échanges d'individus entre les
                populations séparées par la route et les milieux.
          \item[Hypothèse 2.2 :] Les milieux humides qui sont liés des deux côtés de
                la route par un ponceau de plus de 3000~mm ont des populations plus
                stables que celles des milieux humides reliés des deux côtés de la
                route par des ponceaux de moins de 3000~mm.
        \end{description}

  \item Bonifier la qualité des habitats des RFGO et des tortues avec des
        recommandations aidant à une meilleure conception des passages fauniques
        pour ce type de faune (emplacement, protocole de construction, suivi).
\end{enumerate}

À la lumière des résultats obtenus :

\begin{itemize}
  \item L'équipe de recherche pourra indiquer dans quelle mesure les différentes
        caractéristiques de conception des passages fauniques contribuent au
        maintien de la connectivité entre les populations de la petite faune et
        particulièrement de l'herpétofaune, séparée par la route, et à réduire la
        mortalité sur la route.
  \item Les résultats permettront de proposer des modifications de conception qui
        favoriseront l'utilisation des passages par la petite faune.
  \item Les recommandations tiendront compte de contraintes existantes aux sites,
        incluant le contexte hivernal et l'entretien normé des structures, afin de
        permettre aux passages de demeurer fonctionnels pendant leur durée de vie
        utile.
\end{itemize}

Dans le cadre de ce projet, ce sont plus spécifiquement les passages fauniques dans
la région métropolitaine de Montréal favorisant la connectivité pour les RFGO et
pour les tortues qui sont visés. Pour pouvoir répondre à ces objectifs, l'équipe de
recherche a élargi l'aire d'étude à la région métropolitaine de Montréal, pour ne
plus seulement se concentrer sur les sections de l'A10 et de la R104 visées pour des
travaux du MTMD dans le devis initial.

L'échéancier présenté reste inchangé. La suite du rapport présente l'échantillonnage
qui a débuté avec la campagne de terrain 2025, et qui se poursuivra en 2026.
L'année 2027 correspond à la remise du rapport final.

\section{Aire d'étude}

Le projet de recherche tel que révisé dans l'avenant se déroule sur un total de 60
tronçons de 1~km le long des axes principaux du réseau routier du Montréal
métropolitain. Sur la rive sud, nous avons sélectionné des tronçons sur les
autoroutes A10, A15, A20, A30, A40 et A730, et sur les routes secondaires R104,
R112, R205 et R220. Sur la rive nord, nous avons sélectionné des tronçons sur les
autoroutes A15, A20, A30 et A640, et sur les routes secondaires R158 et R344. Chaque
tronçon sera visité au moins deux fois au cours du projet (environ 30 tronçons par
année). Les tronçons de la rive sud traversent l'habitat de différentes espèces à
statut de conservation, notamment la RFGO et la tortue géographique
(\textit{Graptemys geographica}). De ces tronçons, une dizaine ont été identifiés
comme étant prioritaires pour la conservation de la RFGO et de la tortue
géographique, et font l'objet d'un suivi plus fréquent (au moins une fois par
année).

Chaque tronçon inclus dans l'étude contient au moins un ponceau M012 (700~mm à
3000~mm d'ouverture) ou GSQ (plus de 3000~mm d'ouverture). Les sites avec des
ponceaux sont intéressants puisqu'ils signifient qu'il y a présence de milieux
hydriques ou humides à proximité. Nous avons essayé d'inclure le plus de ponceaux
possible dans un même tronçon d'un kilomètre. Le nombre de ponceaux par tronçon
varie donc de 1 à 5. Nous avons priorisé les ponceaux M012 et GSQ qui avaient été
présélectionnés par l'équipe du ministère des Transports et de la Mobilité durable
(MTMD). Ces deux couches de données étaient des sélections spatiales qui
comprenaient des milieux humides de part et d'autre de la route avec un cours d'eau
dans un rayon de 500~m du ponceau. Nous avons également priorisé les tronçons qui
étaient proches de l'habitat connu de la RFGO dans la région métropolitaine de
Montréal. De plus, un site de compensation pour la perte de milieu humide a été
identifié en bordure de l'A30 sur un lot appartenant au Ministère pour une
restauration à venir et une amélioration de l'habitat pour la RFGO. Deux tronçons
ont donc été choisis dans ce secteur, en plus de garder trois tronçons aux sites de
l'A10 et de la R104 pour assurer un suivi avec l'année d'inventaire 2024.

Nous voulons comparer trois catégories de tronçons différents afin d'évaluer
l'influence des milieux à proximité des axes routiers sur la mortalité routière, en
nous concentrant sur la petite faune. Les trois catégories déterminées sont les
suivantes :

\begin{itemize}
  \item \textbf{Catégorie 0} : tronçon témoin sans milieux humides ou zones
        agricoles à moins de 500~m des axes routiers à l'étude ;
  \item \textbf{Catégorie 1} : tronçon avec un ou plusieurs milieux d'intérêt
        présents entre 150 et 300~m du ou des ponceaux d'un seul côté de l'axe
        routier ;
  \item \textbf{Catégorie 2} : tronçon avec un ou plusieurs milieux d'intérêt
        présents à moins de 150~m du ou des ponceaux des deux côtés de l'axe
        routier.
\end{itemize}

Les milieux d'intérêt pour la petite faune que nous avons pris en compte sont les
suivants :

\begin{itemize}
  \item milieux humides d'intérêt de la Communauté métropolitaine de Montréal
        (CMM) ;
  \item milieux humides potentiels du ministère de l'Environnement, de la Lutte
        contre les changements climatiques, de la Faune et des Parcs ;
  \item milieux humides du MTMD ;
  \item habitat désigné de la RFGO par la CMM ;
  \item zonage agricole de la Commission de protection du territoire agricole du
        Québec.
\end{itemize}

Afin d'obtenir une meilleure robustesse statistique lors de la phase d'analyse, nous
avons sélectionné 20 tronçons dans chaque catégorie, ce qui fait un total de 60
sites d'un kilomètre. Nous avons également tenu compte des accès aux tronçons, pour
qu'ils soient les plus sécuritaires possible. De plus, nous avons éliminé des
portions des autoroutes et des routes comprenant des travaux routiers prévus sur le
réseau entre 2025 et 2026. Le processus cartographique complet est disponible dans
le rapport de Fayard et Mazerolle, rendu au MTMD en décembre 2024.

Une visite de repérage des 30 tronçons sélectionnés pour l'année 2025 a été faite
entre le 23 avril et le 2 mai 2025. Cette première phase de terrain a permis de :

\begin{itemize}
  \item visiter les ponceaux qui se trouvaient dans les tronçons routiers afin de
        localiser les meilleurs emplacements pour poser un piège photographique ;
  \item valider les emplacements de stations à reptiles et leur pertinence en
        fonction du milieu ;
  \item vérifier si les vols de drone étaient autorisés ou non dans ce secteur
        (restrictions à proximité des aéroports) ;
  \item repérer des milieux humides proches des ponceaux qui pourraient être suivis
        durant ce projet pour observer l'herpétofaune.
\end{itemize}

Une fois ce repérage terminé, les 30 tronçons routiers inventoriés pour 2025 ont été
validés par le MTMD, 10 de chacune des trois catégories énumérées précédemment. En
tout, nous avons 9 tronçons sur la rive nord et 21 tronçons sur la rive sud de l'île
de Montréal. Le vol de drone est autorisé sur 20 des 30 tronçons.

Une seconde sélection devra avoir lieu pour les sites de 2026. Dix tronçons témoins
de catégorie 0 et dix tronçons avec un milieu humide d'un seul côté de catégorie 1
devront être validés par l'équipe de recherche et le MTMD après repérage.
